\chapter{Introduction}

Financial products have been around for centuries,
mostly in form of credit and insurance,
more recently in form of shared ownership of companies.
This products gave born to their natural extension, the derivatives,
contracts whose value depends on the value of an underlying asset,
they offer ways to transfer risk, also known as hedging,
and to speculate on the future of the underlying,
without owning it.
For many knowing the price to buy or sell such contracts is crucial for their business,
thus the more popular and more traded a product is,
the more interest of modeling its price.
This brings us to 1973 when coincidentally two big events for financial options,
happened.

An option is a contract that gives its holder the right,
but not the obligation,
to buy or sell an asset at an agreed price on an agreed date.
In 1973 Black and Scholes \cite{Black1973} published their now famous model to price European options,
Merton \cite{Merton1973} published the same result derived from a different and more general perspective.
This model set the foundations to price options that coincidentally started to be listed in the CBOE,
a large financial exchange in the US.

As we will see in Chapter \ref{ch:bs},
Black-Scholes model through some assumptions on how the underlying asset dynamics are,
models the price of an option as a partial differential equation (PDE),
\begin{equation*}
    \frac{\partial f}{\partial t} =
        r f(t, S_t)
        - r \frac{\partial f}{\partial S} S_t 
        - \frac{1}{2} \frac{\partial^2 f}{\partial S^2} \sigma^2 S_t^2,
\end{equation*}
one assumption is that the volatility $\sigma$ is constant,
meaning that the variance of the underlying asset dynamics is constant in time,
such simplification allows the PDE to have a closed solution,
\begin{equation*}
    f(t, S_t) = S_t \Phi(d_1) - K e^{-r(T-t)} \Phi(d_2),
\end{equation*}
making the model very popular and useful in practice.

Before 1987,
the empirical evidence backed the constant volatility hypothesis,
but after the 1987 crash,
markets start displaying what is known as the smile,
volatility was not constant anymore,
it depended on the strike price of the option.
The change made the Black-Scholes model inconsistent with the market,
and woke up the interest of modeling the volatility as a function of strike and time.
Many models have been proposed to model the volatility,
we will focus on the one proposed by Dupire \cite{dupire_pricing_smile},
known as the local volatility model,
where the volatility is a deterministic function of time and spot price,
\begin{equation*}
    \sigma(t, S)^2 = 2\,\frac{
        \dfrac{\partial C}{\partial T} + qC + (r-q)\,K\,\dfrac{\partial C}{\partial K}
    }{
        K^2\,\dfrac{\partial^2 C}{\partial K^2}
    }\Bigg|_{\substack{K = S,\\ T = t}},
\end{equation*}
and also models the price of an option as a PDE,
\begin{equation*}
    \frac{\partial C}{\partial T} + (r-q)\,K\,\frac{\partial C}{\partial K}
    - \frac{\sigma^2(T,K)}{2}\,K^2\,\frac{\partial^2 C}{\partial K^2} = -qC.
\end{equation*}
We will see it more in detail in Chapter \ref{ch:local_vol}.

In practice,
all this models require a critical step to be useful,
knowing the value of their parameters,
in particular the volatility.
Local volatility model requires to know the function describing the local volatility surface.
In PDE theory this is known as an inverse problem,
where the goal is to find the parameters of a PDE given some observations of its solution.

\section{Objectives}

The main goal of this thesis is to present a method to calibrate the local volatility model,
while also presenting the theory behind the pricing model,
going through some basics of stochastic calculus, the Black-Scholes equation and arbitrage theory.
So will set a base to understand the local volatility model and its calibration as an inverse problem.

We also aim to do a numerical implementation of both Black-Scholes and local volatility models,
and their calibration,
to check the theory, compare them and test some of the results presented in the literature.
All the code can be found in the repository \href{https://github.com/M315/TFM}{https://github.com/M315/TFM}.

\section{Structure of the document}

Chapters 2 and 3 collect the preliminary notions on probability spaces,
filtrations, martingales and Brownian motion,
the It\^{o} integral, the It\^{o} formula and stochastic differential equations.

Chapter 4 derives the Black-Scholes model,
solves its PDE to obtain the pricing formula,
and defines the implied volatility and its surface.
There is numerical implementation of the implied volatility surface with a more detailed description in Appendix \ref{app:impl_vol}.

Chapter 5 generalizes the pricing argument,
presenting the PDE and the martingale methods for a general contingent claim and the two fundamental theorems of asset pricing.

Chapter 6 introduces the local volatility model,
deriving the Dupire formula and the forward equation,
states the no-arbitrage conditions on the price surface,
and collects some results on the local volatility model.

Chapter \ref{cha:cal_local_vol} treats the calibration problem.
It explains the ill-posedness,
sets up the parameter-to-solution map and the Tikhonov functional,
computes the gradient by the adjoint method,
reviews the convergence rates and the source condition,
and describes the discretization and the minimization.
Here again we present a numerical implementation of the calibration method with a more detailed description in Appendix \ref{app:local_vol}.
